
En
effet, on prend $m=P_{r-1}(c_0,\ldots,c_{r-1})$ avec $c_i=b_i-a_i$, et $M=\sum a_i\binom{m+i}{i}$, de sorte qu'on a une suite exacte
\[
   0\longrightarrow I\longrightarrow {O}_P^{\oplus M}
   \longrightarrow F(m)\longrightarrow0
\]
en vertu de 1.11 et 1.3. De m\^eme, tenant compte de 6.2, on prend $m_1=m+p$ avec $p=P_r(a_{m,0},\ldots,a_{m,r})$, et $M_1=M\cdot\binom{p+N}{N}-\sum a_i\binom{m_1+i}{i}\quad(=\chi(I(p)))$.

\textbf{Th\'eor\`eme 6.4.} \emph{Sous les conditions de 6.1, on suppose $F$ un quotient de ${O}_P(-m)^{\oplus M}$ avec $m\ge0$. Soit $b_i=P_{r-i}(c_{m,i},\ldots,c_{m,r})$ pour $i=0,\ldots,r$, o\`u $P_s$ est le $s$-i\`eme $(b)$-polyn\^ome (1.9)\,; soit $b=b_1+m-1$, et soit $B=\sum_{i=1}^{r} a_i\binom{b+i}{i}$. Alors}~:
\begin{enumerate}[label=(\roman*)]
\item \emph{$F$ est $b$-r\'egulier.}
\item \emph{$-B\le a_0=h^0(F(b))-B$}
\item \emph{$F$ est un $(b_0,\ldots,b_r)$-faisceau.}
\end{enumerate}

En effet, prenons une section $\sigma$ de ${O}_P(1)$ telle qu'on obtienne un diagramme commutatif avec des lignes et colonnes exactes
\[
\begin{tikzcd}[column sep=large,row sep=large]
&0 \arrow[d]&0 \arrow[d]&0 \arrow[d]&\\
0 \arrow[r]&I(-1) \arrow[r] \arrow[d]&I \arrow[r] \arrow[d]&J \arrow[r] \arrow[d]&0\\
0 \arrow[r]&{O}_P(-1)^{\oplus M} \arrow[r] \arrow[d]&{O}_P^{\oplus M} \arrow[r] \arrow[d]&{O}_Y^{\oplus M} \arrow[r] \arrow[d]&0\\
0 \arrow[r]&F(m-1) \arrow[r] \arrow[d]&F(m) \arrow[r] \arrow[d]&G(m) \arrow[r] \arrow[d]&0\\
&0&0&0&.
\end{tikzcd}
\]
